% =============================================================================
% Rapport de Projet : Acc\'{e}l\'{e}rateur Mat\'{e}riel SHA-256 en VHDL
% =============================================================================
\documentclass[12pt,a4paper]{article}

% --- Packages ---
\usepackage[utf8]{inputenc}
\usepackage[T1]{fontenc}
\usepackage[french]{babel}
\usepackage{geometry}
\geometry{margin=2.5cm}
\usepackage{graphicx}
\usepackage{tikz}
\usetikzlibrary{shapes,arrows.meta,positioning,fit,calc,decorations.pathreplacing}
\usepackage{listings}
\usepackage{xcolor}
\usepackage{booktabs}
\usepackage{caption}
\usepackage{subcaption}
\usepackage{amsmath,amssymb}
\usepackage{hyperref}
\usepackage{float}
\usepackage{enumitem}
\usepackage{array}
\usepackage{fancyhdr}
\usepackage{titlesec}

% --- Page style ---
\pagestyle{fancy}
\fancyhf{}
\fancyhead[L]{\footnotesize Acc\'{e}l\'{e}rateur SHA-256 -- DIMRI Imad}
\fancyhead[R]{\footnotesize \thepage}
\renewcommand{\headrulewidth}{0.4pt}

% --- Section formatting ---
\titleformat{\section}{\Large\bfseries}{\thesection}{1em}{}
\titleformat{\subsection}{\large\bfseries}{\thesubsection}{1em}{}

% --- Code listing style ---
\definecolor{codebg}{RGB}{248,248,248}
\definecolor{codegreen}{RGB}{0,128,0}
\definecolor{codegray}{RGB}{100,100,100}
\definecolor{codeblue}{RGB}{0,0,160}
\definecolor{codered}{RGB}{160,0,0}
\lstdefinestyle{vhdlstyle}{
    language=VHDL,
    backgroundcolor=\color{codebg},
    basicstyle=\ttfamily\scriptsize,
    keywordstyle=\color{codeblue}\bfseries,
    commentstyle=\color{codegreen}\itshape,
    stringstyle=\color{codered},
    numbers=left,
    numberstyle=\tiny\color{codegray},
    numbersep=5pt,
    frame=single,
    breaklines=true,
    captionpos=b,
    tabsize=4,
    xleftmargin=1em,
    framexleftmargin=1em
}
\lstset{style=vhdlstyle}

% --- Title ---
\title{%
\vspace{-1cm}
\rule{\textwidth}{1.5pt}\\[0.4cm]
{\Huge\bfseries Conception et Validation d'un\\Acc\'{e}l\'{e}rateur Mat\'{e}riel SHA-256}\\[0.3cm]
{\Large Co-processeur cryptographique synth\'{e}tisable sur FPGA\\
avec support multi-blocs}\\[0.2cm]
\rule{\textwidth}{1.5pt}
}
\author{%
\textbf{DIMRI Imad}\\[0.2cm]
\small Projet de conception num\'{e}rique
}
\date{Mai 2026}

\begin{document}
\maketitle
\thispagestyle{empty}

\begin{abstract}
\noindent
Ce rapport pr\'{e}sente la conception, l'impl\'{e}mentation et la v\'{e}rification
fonctionnelle d'un \textbf{acc\'{e}l\'{e}rateur mat\'{e}riel} d\'{e}di\'{e} au calcul de la
fonction de hachage cryptographique SHA-256. Impl\'{e}ment\'{e} en VHDL et valid\'{e}
sous Vivado XSim (AMD/Xilinx), le co-processeur supporte le hachage de
messages de longueur arbitraire via un traitement multi-blocs. L'architecture
it\'{e}rative \`{a} une ronde par cycle d'horloge permet d'atteindre un d\'{e}bit
th\'{e}orique de \textbf{602 Mbit/s} \`{a} 100 MHz, soit un gain de performance de
\textbf{24\`{a} 48 fois} par rapport \`{a} une ex\'{e}cution logicielle sur processeur
embarqu\'{e}. La conformit\'{e} fonctionnelle est d\'{e}montr\'{e}e par simulation avec les
vecteurs de test officiels NIST, incluant un message n\'{e}cessitant deux blocs
de 512 bits.
\end{abstract}

\vspace{0.5cm}
\tableofcontents
\newpage


% =============================================================================
\section{Introduction}
\label{sec:intro}
% =============================================================================

\subsection{Contexte et motivation}

Les fonctions de hachage cryptographique sont au c\oe{}ur de la s\'{e}curit\'{e}
num\'{e}rique moderne. SHA-256 (\emph{Secure Hash Algorithm}, 256 bits), d\'{e}fini
par la norme FIPS 180-4 du NIST~\cite{fips180}, est utilis\'{e} dans :
\begin{itemize}[leftmargin=*]
    \item La \textbf{signature \'{e}lectronique} (certificats X.509, TLS/SSL) ;
    \item La \textbf{v\'{e}rification d'int\'{e}grit\'{e}} (checksums, mises \`{a} jour firmware) ;
    \item Les \textbf{protocoles blockchain} (Bitcoin, Ethereum PoW) ;
    \item L'\textbf{authentification} (HMAC-SHA256, d\'{e}rivation de cl\'{e}s).
\end{itemize}

Sur un processeur embarqu\'{e} (ARM Cortex-M4 par exemple), le calcul de SHA-256
n\'{e}cessite environ 3\,500 cycles par bloc de 512 bits, ce qui constitue un
goulot d'\'{e}tranglement pour les applications \`{a} haut d\'{e}bit. Un
\textbf{acc\'{e}l\'{e}rateur mat\'{e}riel} d\'{e}di\'{e} permet de r\'{e}duire ce co\^{u}t \`{a}
$\sim$85 cycles par bloc, soit un gain de plus de 40 fois.

\subsection{Objectifs du projet}

\begin{enumerate}
    \item Impl\'{e}menter l'algorithme SHA-256 conform\'{e}ment \`{a} FIPS 180-4 ;
    \item Supporter les messages de longueur arbitraire (multi-blocs) ;
    \item Adopter une architecture modulaire, synth\'{e}tisable sur FPGA ;
    \item Valider par simulation avec les vecteurs de test NIST ;
    \item Analyser les performances et justifier le terme \og acc\'{e}l\'{e}rateur \fg{}.
\end{enumerate}

\subsection{Organisation du rapport}

\begin{itemize}
    \item Section~\ref{sec:algo} : rappel de l'algorithme SHA-256 ;
    \item Section~\ref{sec:archi} : architecture mat\'{e}rielle du co-processeur ;
    \item Section~\ref{sec:impl} : impl\'{e}mentation VHDL module par module ;
    \item Section~\ref{sec:verif} : v\'{e}rification fonctionnelle et r\'{e}sultats ;
    \item Section~\ref{sec:accel} : analyse de performance et justification du terme acc\'{e}l\'{e}rateur ;
    \item Section~\ref{sec:bugs} : probl\`{e}mes rencontr\'{e}s et corrections ;
    \item Section~\ref{sec:conclu} : conclusion et perspectives.
\end{itemize}

% =============================================================================
\section{Algorithme SHA-256}
\label{sec:algo}
% =============================================================================

\subsection{Vue d'ensemble}

SHA-256 transforme un message de longueur arbitraire en un condens\'{e} fixe de
256 bits (32 octets). Le traitement se d\'{e}compose en trois phases principales,
illustr\'{e}es par la figure~\ref{fig:algo_overview}.

\begin{figure}[H]
\centering
% =====================================================================
% IMAGE PLACEHOLDER: Vue d'ensemble de l'algorithme SHA-256
% Ins\'{e}rer ici un sch\'{e}ma montrant les 3 phases :
%   1. Padding du message
%   2. D\'{e}coupage en blocs de 512 bits
%   3. Compression it\'{e}rative (64 rondes par bloc)
% =====================================================================
\fbox{\parbox{0.85\textwidth}{\centering\vspace{3cm}
\textbf{[IMAGE : Vue d'ensemble de l'algorithme SHA-256]}\\[0.3cm]
\small Sch\'{e}ma montrant : Message $\rightarrow$ Padding $\rightarrow$
Blocs 512 bits $\rightarrow$ 64 rondes de compression $\rightarrow$ Digest 256 bits
\vspace{3cm}}}
\caption{Vue d'ensemble du traitement SHA-256}
\label{fig:algo_overview}
\end{figure}

\subsection{Phase 1 : Pr\'{e}-traitement (Padding)}

Le message est compl\'{e}t\'{e} pour que sa longueur soit un multiple de 512 bits :
\begin{enumerate}
    \item Ajout d'un bit \texttt{'1'} imm\'{e}diatement apr\`{e}s le message ;
    \item Ajout de bits \texttt{'0'} jusqu'\`{a} ce que la longueur soit
          congrue \`{a} 448 modulo 512 ;
    \item Ajout de la longueur originale du message sur 64 bits (big-endian).
\end{enumerate}

\textbf{Exemple} : pour le message \texttt{"abc"} (24 bits) :
\begin{verbatim}
    61626380 00000000 00000000 00000000
    00000000 00000000 00000000 00000000
    00000000 00000000 00000000 00000000
    00000000 00000000 00000000 00000018
\end{verbatim}
Le r\'{e}sultat est un bloc unique de 512 bits (16 mots de 32 bits).

\subsection{Phase 2 : Expansion du message}

Les 16 mots d'entr\'{e}e ($W_0 \ldots W_{15}$) sont \'{e}tendus en 64 mots par la
r\'{e}currence :
\begin{equation}
W_t = \sigma_1(W_{t-2}) + W_{t-7} + \sigma_0(W_{t-15}) + W_{t-16}
\quad \text{pour } 16 \leq t \leq 63
\label{eq:expansion}
\end{equation}

\subsection{Phase 3 : Compression (64 rondes)}

Chaque ronde $t$ calcule deux valeurs temporaires (tableau~\ref{tab:fonctions}) :
\begin{align}
T_1 &= h + \Sigma_1(e) + Ch(e,f,g) + K_t + W_t \label{eq:t1}\\
T_2 &= \Sigma_0(a) + Maj(a,b,c) \label{eq:t2}
\end{align}

Puis les huit variables de travail sont mises \`{a} jour :
\[
h \leftarrow g,\quad g \leftarrow f,\quad f \leftarrow e,\quad
e \leftarrow d + T_1,\quad d \leftarrow c,\quad c \leftarrow b,\quad
b \leftarrow a,\quad a \leftarrow T_1 + T_2
\]

\begin{table}[H]
\centering
\caption{Fonctions logiques utilis\'{e}es dans SHA-256}
\label{tab:fonctions}
\begin{tabular}{ll}
\toprule
\textbf{Fonction} & \textbf{D\'{e}finition (op\'{e}rations sur 32 bits)} \\
\midrule
$Ch(x,y,z)$        & $(x \wedge y) \oplus (\neg x \wedge z)$ \\
$Maj(x,y,z)$       & $(x \wedge y) \oplus (x \wedge z) \oplus (y \wedge z)$ \\
$\Sigma_0(x)$      & $ROTR^{2}(x) \oplus ROTR^{13}(x) \oplus ROTR^{22}(x)$ \\
$\Sigma_1(x)$      & $ROTR^{6}(x) \oplus ROTR^{11}(x) \oplus ROTR^{25}(x)$ \\
$\sigma_0(x)$      & $ROTR^{7}(x) \oplus ROTR^{18}(x) \oplus SHR^{3}(x)$ \\
$\sigma_1(x)$      & $ROTR^{17}(x) \oplus ROTR^{19}(x) \oplus SHR^{10}(x)$ \\
\bottomrule
\end{tabular}
\end{table}

\subsection{Accumulation multi-blocs}

Pour un message compos\'{e} de $N$ blocs, le condens\'{e} interm\'{e}diaire est accumul\'{e}
apr\`{e}s chaque bloc :
\[
H^{(i)}_j = H^{(i-1)}_j + \text{variable}_j \quad \text{pour } j = 0 \ldots 7
\]
avec $H^{(0)} = H_{init}$ (les huit constantes d'initialisation).


% =============================================================================
\section{Architecture mat\'{e}rielle}
\label{sec:archi}
% =============================================================================

\subsection{Choix architectural}

L'architecture adopt\'{e}e est de type \textbf{it\'{e}rative} : une seule instance
du circuit de ronde est instanci\'{e}e et r\'{e}utilis\'{e}e 64 fois. Ce choix offre
le meilleur compromis entre surface et d\'{e}bit pour un co-processeur embarqu\'{e} :

\begin{itemize}
    \item Surface mod\'{e}r\'{e}e (un seul jeu d'additionneurs 32 bits) ;
    \item D\'{e}bit \'{e}lev\'{e} ($\sim$85 cycles/bloc = 602 Mbit/s \`{a} 100 MHz) ;
    \item Complexit\'{e} de contr\^{o}le limit\'{e}e (une FSM lin\'{e}aire).
\end{itemize}

\subsection{Diagramme de blocs global}

La figure~\ref{fig:block_diagram} pr\'{e}sente l'architecture compl\`{e}te du
co-processeur avec ses interconnexions.

\begin{figure}[H]
\centering
% =====================================================================
% IMAGE PLACEHOLDER: Diagramme de blocs du co-processeur
% Ins\'{e}rer ici le diagramme montrant :
%   - SHA256_FSM (contr\^{o}le) au centre
%   - Scheduler (expansion W) \`{a} droite
%   - Round (compression) en bas \`{a} droite
%   - Registres a..h + H en bas
%   - M\'{e}moire partag\'{e}e \`{a} gauche
%   - Fl\`{e}ches montrant les signaux de contr\^{o}le et de donn\'{e}es
% =====================================================================
\fbox{\parbox{0.85\textwidth}{\centering\vspace{4cm}
\textbf{[IMAGE : Diagramme de blocs du co-processeur SHA-256]}\\[0.3cm]
\small Modules : FSM, Scheduler, Round, Registres (a..h + H), M\'{e}moire.\\
Signaux : start/done, addr/data, $W_t$, $K_t$, round\_en, hash\_init, compute\_done.
\vspace{4cm}}}
\caption{Architecture globale du co-processeur SHA-256}
\label{fig:block_diagram}
\end{figure}

\subsection{Description des modules}

Le syst\`{e}me est d\'{e}compos\'{e} en six modules VHDL, r\'{e}sum\'{e}s dans le
tableau~\ref{tab:modules}.

\begin{table}[H]
\centering
\caption{Modules VHDL du co-processeur et leurs r\^{o}les}
\label{tab:modules}
\begin{tabular}{lp{9cm}}
\toprule
\textbf{Module} & \textbf{R\^{o}le} \\
\midrule
\texttt{SHA256\_pkg}    & Package contenant les d\'{e}finitions de types (\texttt{word},
                          tableaux), les 64 constantes $K_t$, les valeurs initiales
                          $H_{init}$, et toutes les fonctions logiques. \\[4pt]
\texttt{Round}          & Circuit \textbf{purement combinatoire} calculant une ronde
                          de compression ($T_1$, $T_2$, nouvelles variables). \\[4pt]
\texttt{Scheduler}      & Fen\^{e}tre glissante de 16 registres impl\'{e}mentant
                          l'expansion du message (r\'{e}currence~\ref{eq:expansion}). \\[4pt]
\texttt{SHA256\_FSM}    & Machine \`{a} \'{e}tats finis g\'{e}rant le s\'{e}quencement,
                          le pipeline d'adresses m\'{e}moire, la s\'{e}lection de $K_t$,
                          et la boucle multi-blocs. \\[4pt]
\texttt{SHA256\_Core}   & Entit\'{e} de plus haut niveau : instancie les modules
                          ci-dessus et contient les registres de travail $a \ldots h$
                          et l'accumulateur $H$. \\[4pt]
\texttt{SHA256\_Memory} & RAM synchrone \`{a} double port (16 Ko), latence de lecture
                          d'un cycle d'horloge. \\
\bottomrule
\end{tabular}
\end{table}

\subsection{Carte m\'{e}moire}

L'interface entre le processeur h\^{o}te et le co-processeur utilise une m\'{e}moire
partag\'{e}e dont l'organisation est d\'{e}crite dans le tableau~\ref{tab:memmap}.

\begin{table}[H]
\centering
\caption{Organisation de la m\'{e}moire partag\'{e}e}
\label{tab:memmap}
\begin{tabular}{lll}
\toprule
\textbf{Adresse (octet)} & \textbf{Taille} & \textbf{Contenu} \\
\midrule
\texttt{0x0000}                    & 1 mot   & Nombre de blocs $N$ \\
\texttt{0x0004} + $b{\times}64$ + $w{\times}4$ & $N{\times}16$ mots & Donn\'{e}es pr\'{e}-padd\'{e}es \\
\texttt{0x2000}--\texttt{0x201C}   & 8 mots  & Condens\'{e} r\'{e}sultant $H_0 \ldots H_7$ \\
\bottomrule
\end{tabular}
\end{table}

Le h\^{o}te effectue le padding en logiciel avant d'\'{e}crire les blocs en m\'{e}moire.
Le co-processeur se concentre sur le calcul intensif --- c'est cette s\'{e}paration
des responsabilit\'{e}s qui caract\'{e}rise un acc\'{e}l\'{e}rateur mat\'{e}riel.

\subsection{Protocole d'utilisation}

\begin{enumerate}
    \item Le h\^{o}te \'{e}crit le nombre de blocs \`{a} l'adresse \texttt{0x0000} ;
    \item Le h\^{o}te \'{e}crit les blocs pr\'{e}-padd\'{e}s \`{a} partir de \texttt{0x0004} ;
    \item Le h\^{o}te pulse \texttt{start} pendant un cycle (\texttt{done = '1'}) ;
    \item Le co-processeur traite tous les blocs automatiquement ;
    \item Quand \texttt{done} repasse \`{a} '1', le r\'{e}sultat est disponible \`{a}
          \texttt{0x2000}.
\end{enumerate}

\subsection{Machine \`{a} \'{e}tats finis (FSM)}

La FSM comporte sept \'{e}tats, illustr\'{e}s par la figure~\ref{fig:fsm_diagram}.

\begin{figure}[H]
\centering
% =====================================================================
% IMAGE PLACEHOLDER: Diagramme d'\'{e}tats de la FSM
% Ins\'{e}rer ici le diagramme d'\'{e}tats montrant :
%   IDLE -> READ_NBLK -> LOAD (16 cyc) -> COMPUTE (49 cyc) -> ACCUM
%   ACCUM -> PREFETCH -> LOAD (boucle bloc suivant)
%   ACCUM -> WRITEBACK -> IDLE (dernier bloc)
% =====================================================================
\fbox{\parbox{0.85\textwidth}{\centering\vspace{3.5cm}
\textbf{[IMAGE : Diagramme d'\'{e}tats de la FSM multi-blocs]}\\[0.3cm]
\small \'{E}tats : IDLE, READ\_NBLK, PREFETCH, LOAD, COMPUTE, ACCUM, WRITEBACK.\\
Transitions : start=1, counter=15, counter=64, bloc\_suivant, dernier\_bloc, wb\_cnt=7.
\vspace{3.5cm}}}
\caption{Diagramme d'\'{e}tats de la FSM de contr\^{o}le}
\label{fig:fsm_diagram}
\end{figure}

\noindent
\textbf{Description des \'{e}tats :}
\begin{itemize}[leftmargin=*]
    \item \textbf{IDLE} : attente du signal \texttt{start} ; adresse 0x0000
          maintenue sur le bus pour pr\'{e}-lecture du nombre de blocs.
    \item \textbf{READ\_NBLK} : capture du nombre de blocs ; initialisation de
          $H$ \`{a} $H_{init}$ ; copie de $H$ dans les registres de travail.
    \item \textbf{PREFETCH} : utilis\'{e} pour les blocs suivants ; pr\'{e}pare le
          pipeline d'adresses et copie le $H$ accumul\'{e} dans les registres.
    \item \textbf{LOAD} : 16 cycles chargeant les mots $W_0 \ldots W_{15}$ dans
          le Scheduler ; les rondes 0--14 s'ex\'{e}cutent en parall\`{e}le.
    \item \textbf{COMPUTE} : 49 cycles ex\'{e}cutant les rondes 15--63.
    \item \textbf{ACCUM} : accumulation $H^{(i)} = H^{(i-1)} + (a \ldots h)$ ;
          branchement vers PREFETCH (bloc suivant) ou WRITEBACK (dernier).
    \item \textbf{WRITEBACK} : \'{e}criture s\'{e}quentielle des 8 mots du condens\'{e}
          \`{a} l'adresse \texttt{0x2000}.
\end{itemize}


% =============================================================================
\section{Impl\'{e}mentation VHDL}
\label{sec:impl}
% =============================================================================

Cette section pr\'{e}sente les extraits de code essentiels de chaque module.
Le code complet est disponible dans le d\'{e}p\^{o}t du projet.

\subsection{Package : types et fonctions logiques}

Le listing~\ref{lst:pkg} montre les fonctions de rotation et d'expansion
$\sigma_0$. Notons que les rotations (\texttt{ROTR}) sont impl\'{e}ment\'{e}es par
simple concat\'{e}nation --- en mat\'{e}riel, cela ne co\^{u}te \textbf{aucune} ressource
logique (c'est un simple reroutage de fils).

\begin{lstlisting}[caption={Extrait de \texttt{SHA256\_pkg.vhd} -- fonctions logiques},
                   label={lst:pkg}]
function ROTR (x : in word; n : in integer) return word is
begin
    return x(n-1 downto 0) & x(31 downto n);
    -- Zero cost: just wire routing in hardware!
end function ROTR;

function SHR (x : in word; n : in integer) return word is
begin
    return (n-1 downto 0 => '0') & x(31 downto n);
end function SHR;

function small_sigma0 (x : in word) return word is
begin
    return ROTR(x,7) xor ROTR(x,18) xor SHR(x,3);
end function small_sigma0;
\end{lstlisting}

\subsection{Module Round : chemin de donn\'{e}es critique}

Le listing~\ref{lst:round} montre le calcul combinatoire d'une ronde compl\`{e}te.
La figure~\ref{fig:round_critical} illustre le chemin critique.

\begin{figure}[H]
\centering
% =====================================================================
% IMAGE PLACEHOLDER: Chemin critique de la ronde de compression
% Ins\'{e}rer ici un sch\'{e}ma montrant :
%   - Les entr\'{e}es h, e, f, g, Kt, Wt, a, b, c, d
%   - Les blocs fonctionnels : Sigma1, Ch, Sigma0, Maj
%   - La cha\^{i}ne d'additionneurs pour T1 (5 op\'{e}randes)
%   - La sortie T2 et les calculs de new_a et new_e
%   - Le chemin critique mis en \'{e}vidence (en rouge)
% =====================================================================
\fbox{\parbox{0.85\textwidth}{\centering\vspace{4cm}
\textbf{[IMAGE : Chemin de donn\'{e}es d'une ronde SHA-256]}\\[0.3cm]
\small Montrer : $\Sigma_1(e)$, $Ch(e,f,g)$ en parall\`{e}le avec $\Sigma_0(a)$, $Maj(a,b,c)$.\\
Cha\^{i}ne critique en rouge : $h + \Sigma_1 + Ch + K_t + W_t$ (4 additions s\'{e}rie).\\
Sorties : $a' = T_1 + T_2$,\quad $e' = d + T_1$.
\vspace{4cm}}}
\caption{Chemin de donn\'{e}es combinatoire d'une ronde (chemin critique en rouge)}
\label{fig:round_critical}
\end{figure}

\begin{lstlisting}[caption={Module \texttt{Round.vhd} -- calcul d'une ronde compl\`{e}te},
                   label={lst:round}]
t1 <= std_logic_vector( unsigned(h)
                      + unsigned(big_sigma1(e))
                      + unsigned(ch(e, f, g))
                      + unsigned(kt)
                      + unsigned(wt) );

t2 <= std_logic_vector( unsigned(big_sigma0(a))
                      + unsigned(maj(a, b, c)) );

n_h <= g;  n_g <= f;  n_f <= e;
n_e <= std_logic_vector(unsigned(d) + unsigned(t1));
n_d <= c;  n_c <= b;  n_b <= a;
n_a <= std_logic_vector(unsigned(t1) + unsigned(t2));
\end{lstlisting}

\subsection{Module Scheduler : fen\^{e}tre glissante}

Le Scheduler (listing~\ref{lst:sched}) impl\'{e}mente une fen\^{e}tre glissante de
16 registres. Pour $t < 16$, les mots sont charg\'{e}s depuis la m\'{e}moire. Pour
$t \geq 16$, le nouveau mot est calcul\'{e} par la r\'{e}currence de
l'\'{e}quation~\ref{eq:expansion}.

\begin{lstlisting}[caption={Module \texttt{Scheduler.vhd} -- expansion du message},
                   label={lst:sched}]
new_w <= std_logic_vector(
    unsigned(small_sigma1(w_window(14)))
  + unsigned(w_window(9))
  + unsigned(small_sigma0(w_window(1)))
  + unsigned(w_window(0)));

output <= w_window(15);  -- W[t] for the Round

process(clk)
begin
    if rising_edge(clk) then
        if shift_e = '1' then
            w_window(0 to 14) <= w_window(1 to 15);  -- shift left
        end if;
        if w_src = '0' then
            w_window(15) <= input;  -- load from memory (t < 16)
        else
            w_window(15) <= new_w;  -- computed (t >= 16)
        end if;
    end if;
end process;
\end{lstlisting}

\subsection{Module FSM : pipeline d'adresses et boucle multi-blocs}

Le listing~\ref{lst:fsm} montre la partie critique de la FSM : la gestion du
pipeline d'adresses et la boucle entre blocs.

\begin{lstlisting}[caption={Extrait de \texttt{SHA256\_FSM.vhd} -- boucle multi-blocs},
                   label={lst:fsm}]
when READ_NBLK =>
    -- mem_data = num_blocks (from IDLE's steady addr)
    num_blocks <= to_integer(unsigned(mem_data(7 downto 0)));
    block_idx  <= 0;
    binit_i    <= '1';   -- H_result <= H_init
    addr_reg   <= blk_word_addr(0, 1);  -- queue W[1]
    state      <= LOAD;

when ACCUM =>
    if block_idx + 1 < num_blocks then
        -- Next block: drive W[0] address
        block_idx <= block_idx + 1;
        addr_reg  <= blk_word_addr(block_idx + 1, 0);
        state     <= PREFETCH;
    else
        wb_cnt <= 0;
        state  <= WRITEBACK;  -- done!
    end if;
\end{lstlisting}

\subsection{Module Core : registres et accumulation}

Le listing~\ref{lst:core} montre la logique s\'{e}quentielle g\'{e}rant les registres
de travail et l'accumulation du condens\'{e} entre les blocs.

\begin{lstlisting}[caption={Extrait de \texttt{SHA256\_Core.vhd} -- accumulation multi-blocs},
                   label={lst:core}]
-- Block init: H = H_init (first block only)
if block_init = '1' then
    H_result <= H_init;
end if;

-- Per-block: copy H -> working registers
if hash_init = '1' then
    ra <= H_result(0); rb <= H_result(1);
    rc <= H_result(2); rd <= H_result(3);
    re <= H_result(4); rf <= H_result(5);
    rg <= H_result(6); rh <= H_result(7);
elsif round_en = '1' then
    ra <= n_a; rb <= n_b; rc <= n_c; rd <= n_d;
    re <= n_e; rf <= n_f; rg <= n_g; rh <= n_h;
end if;

-- End of block: accumulate H += working vars
if compute_done = '1' then
    H_result(0) <= add(H_result(0), ra);
    H_result(1) <= add(H_result(1), rb);
    -- ... H_result(2) to H_result(7)
end if;
\end{lstlisting}


% =============================================================================
\section{V\'{e}rification fonctionnelle et r\'{e}sultats}
\label{sec:verif}
% =============================================================================

\subsection{Strat\'{e}gie de v\'{e}rification}

La v\'{e}rification repose sur un \textbf{banc de test auto-v\'{e}rifiant}
(\texttt{SHA256\_tb.vhd}) qui automatise enti\`{e}rement le processus :
\begin{enumerate}
    \item Pr\'{e}-chargement d'un message padd\'{e} dans la m\'{e}moire ;
    \item Activation du co-processeur (\texttt{start}) ;
    \item Capture des 8 mots \'{e}crits pendant WRITEBACK ;
    \item Comparaison mot par mot avec le condens\'{e} de r\'{e}f\'{e}rence.
\end{enumerate}

Les condens\'{e}s de r\'{e}f\'{e}rence sont calcul\'{e}s avec Python (\texttt{hashlib.sha256})
et v\'{e}rifi\'{e}s ind\'{e}pendamment.

\subsection{Vecteurs de test}

Trois vecteurs couvrent les cas critiques du design
(tableau~\ref{tab:vectors}).

\begin{table}[H]
\centering
\caption{Vecteurs de test utilis\'{e}s pour la v\'{e}rification}
\label{tab:vectors}
\begin{tabular}{llcc}
\toprule
\textbf{Message} & \textbf{Longueur} & \textbf{Blocs} & \textbf{Cas test\'{e}} \\
\midrule
\texttt{"abc"} & 24 bits & 1 & Padding non align\'{e} sur mot \\
\texttt{"abcd"} & 32 bits & 1 & Padding align\'{e} sur mot \\
56 octets NIST$^*$ & 448 bits & 2 & \textbf{Multi-blocs} \\
\bottomrule
\end{tabular}\\[4pt]
{\footnotesize $^*$ \texttt{"abcdbcdecdefdefgefghfghighijhijkijkljklmklmnlmnomnopnopq"}}
\end{table}

Le vecteur \`{a} 2 blocs est particuli\`{e}rement important : il v\'{e}rifie que
l'accumulation de $H$ entre les blocs fonctionne correctement.

\subsection{Environnement de simulation}

\begin{itemize}
    \item \textbf{Outil} : AMD Vivado 2025.2, simulateur XSim ;
    \item \textbf{Standard VHDL} : VHDL-93 (aucune d\'{e}pendance 2008) ;
    \item \textbf{P\'{e}riode d'horloge} : 10 ns (100 MHz) ;
    \item \textbf{Dur\'{e}e de simulation} : 5 $\mu$s.
\end{itemize}

\subsection{Capture d'\'{e}cran de simulation}

La figure~\ref{fig:sim_console} montre la sortie console de la simulation
dans Vivado XSim, confirmant le passage de tous les tests.

\begin{figure}[H]
\centering
% =====================================================================
% IMAGE PLACEHOLDER: Capture d'\'{e}cran de la console XSim
% Ins\'{e}rer ici la capture montrant :
%   === Test: abc (1 block) ===
%   H0 ok BA7816BF ... H7 ok F20015AD
%   === Test: abcd (1 block) ===
%   H0 ok 88D4266F ... H7 ok 6F031589
%   === Test: 56-byte NIST (2 blocks) ===
%   H0 ok 248D6A61 ... H7 ok 19DB06C1
%   ALL TESTS PASSED
% =====================================================================
\fbox{\parbox{0.85\textwidth}{\centering\vspace{3.5cm}
\textbf{[IMAGE : Capture de la console XSim -- ALL TESTS PASSED]}\\[0.3cm]
\small Montrer la sortie compl\`{e}te de la simulation avec les 3 vecteurs\\
et le message final \texttt{ALL TESTS PASSED}
\vspace{3.5cm}}}
\caption{Sortie de simulation XSim montrant le passage des 3 vecteurs de test}
\label{fig:sim_console}
\end{figure}

\subsection{Chronogrammes}

La figure~\ref{fig:waveform_overview} montre les signaux principaux pendant
le traitement d'un bloc, et la figure~\ref{fig:waveform_multiblock} montre
la transition entre deux blocs.

\begin{figure}[H]
\centering
% =====================================================================
% IMAGE PLACEHOLDER: Chronogramme d'un bloc complet
% Ins\'{e}rer ici le chronogramme (waveform) montrant sur un bloc :
%   - clk
%   - state (IDLE -> READ_NBLK -> LOAD -> COMPUTE -> ACCUM -> WRITEBACK)
%   - counter (0..15 en LOAD, 16..64 en COMPUTE)
%   - round_en
%   - wt (valeurs changeantes)
%   - mem_wen (pulses en WRITEBACK)
%   - done
% =====================================================================
\fbox{\parbox{0.85\textwidth}{\centering\vspace{3.5cm}
\textbf{[IMAGE : Chronogramme -- traitement d'un bloc complet]}\\[0.3cm]
\small Signaux : clk, state, counter, round\_en, shift\_e, wt, done, mem\_wen.\\
Montrer la progression IDLE $\rightarrow$ READ\_NBLK $\rightarrow$ LOAD $\rightarrow$ COMPUTE $\rightarrow$ ACCUM $\rightarrow$ WRITEBACK.
\vspace{3.5cm}}}
\caption{Chronogramme du traitement d'un bloc SHA-256 complet}
\label{fig:waveform_overview}
\end{figure}

\begin{figure}[H]
\centering
% =====================================================================
% IMAGE PLACEHOLDER: Chronogramme multi-blocs (transition entre blocs)
% Ins\'{e}rer ici le chronogramme montrant la transition entre le bloc 0
% et le bloc 1 pour le test 56 octets :
%   - state passant de ACCUM -> PREFETCH -> LOAD (nouveau bloc)
%   - block_idx changeant de 0 \`{a} 1
%   - H_result se mettant \`{a} jour \`{a} l'ACCUM
% =====================================================================
\fbox{\parbox{0.85\textwidth}{\centering\vspace{3.5cm}
\textbf{[IMAGE : Chronogramme -- transition multi-blocs]}\\[0.3cm]
\small Montrer la transition ACCUM $\rightarrow$ PREFETCH $\rightarrow$ LOAD (bloc 1).\\
Signaux : state, block\_idx, H\_result[0], counter, hash\_init.
\vspace{3.5cm}}}
\caption{Chronogramme montrant la transition entre le bloc 0 et le bloc 1}
\label{fig:waveform_multiblock}
\end{figure}

\subsection{R\'{e}sultats quantitatifs}

Le tableau~\ref{tab:results} r\'{e}sume les temps de compl\'{e}tion mesur\'{e}s en
simulation.

\begin{table}[H]
\centering
\caption{R\'{e}sultats de simulation -- temps et nombre de cycles}
\label{tab:results}
\begin{tabular}{lcccc}
\toprule
\textbf{Test} & \textbf{Blocs} & \textbf{Temps (ns)} & \textbf{Cycles} & \textbf{Statut} \\
\midrule
\texttt{"abc"}         & 1 & 965   & 96  & \textbf{PASS} \\
\texttt{"abcd"}        & 1 & 1\,955 & 99  & \textbf{PASS} \\
56 oct.\ NIST          & 2 & 3\,775 & 188 & \textbf{PASS} \\
\bottomrule
\end{tabular}
\end{table}

\subsection{Condens\'{e}s obtenus vs.\ r\'{e}f\'{e}rence}

Le tableau~\ref{tab:digests} montre les condens\'{e}s complets obtenus pour le
vecteur \texttt{"abc"}, compar\'{e}s \`{a} la valeur officielle NIST.

\begin{table}[H]
\centering
\caption{Condens\'{e} SHA-256 de \texttt{"abc"} -- obtenu vs.\ r\'{e}f\'{e}rence NIST}
\label{tab:digests}
\begin{tabular}{cll}
\toprule
\textbf{Mot} & \textbf{Obtenu (simulation)} & \textbf{R\'{e}f\'{e}rence NIST} \\
\midrule
$H_0$ & \texttt{BA7816BF} & \texttt{BA7816BF} \\
$H_1$ & \texttt{8F01CFEA} & \texttt{8F01CFEA} \\
$H_2$ & \texttt{414140DE} & \texttt{414140DE} \\
$H_3$ & \texttt{5DAE2223} & \texttt{5DAE2223} \\
$H_4$ & \texttt{B00361A3} & \texttt{B00361A3} \\
$H_5$ & \texttt{96177A9C} & \texttt{96177A9C} \\
$H_6$ & \texttt{B410FF61} & \texttt{B410FF61} \\
$H_7$ & \texttt{F20015AD} & \texttt{F20015AD} \\
\bottomrule
\end{tabular}
\end{table}


% =============================================================================
\section{Analyse de performance : pourquoi c'est un acc\'{e}l\'{e}rateur}
\label{sec:accel}
% =============================================================================

\subsection{D\'{e}finition d'un acc\'{e}l\'{e}rateur mat\'{e}riel}

Un \textbf{acc\'{e}l\'{e}rateur mat\'{e}riel} est un circuit sp\'{e}cialis\'{e} qui ex\'{e}cute une
fonction donn\'{e}e significativement plus vite (ou avec moins d'\'{e}nergie) qu'un
processeur g\'{e}n\'{e}raliste. Notre co-processeur SHA-256 satisfait pleinement cette
d\'{e}finition pour trois raisons fondamentales.

\subsection{Raison 1 : Parall\'{e}lisme spatial}

Un processeur g\'{e}n\'{e}raliste ex\'{e}cute les op\'{e}rations \textbf{s\'{e}quentiellement}
(une instruction par cycle dans le meilleur cas). Notre circuit calcule
\emph{simultan\'{e}ment} dans des chemins combinatoires d\'{e}di\'{e}s :
\begin{itemize}
    \item $\Sigma_1(e)$ et $Ch(e,f,g)$ en parall\`{e}le ;
    \item $\Sigma_0(a)$ et $Maj(a,b,c)$ en parall\`{e}le ;
    \item L'expansion $W_t$ (Scheduler) en parall\`{e}le avec la compression.
\end{itemize}

La figure~\ref{fig:parallel_comparison} illustre cette diff\'{e}rence.

\begin{figure}[H]
\centering
% =====================================================================
% IMAGE PLACEHOLDER: Comparaison ex\'{e}cution s\'{e}quentielle vs parall\`{e}le
% Ins\'{e}rer ici un diagramme comparant :
%   - En haut : timeline CPU (instructions s\'{e}quentielles, ~50 par ronde)
%   - En bas : timeline FPGA (tout en 1 cycle, chemins parall\`{e}les)
% =====================================================================
\fbox{\parbox{0.85\textwidth}{\centering\vspace{3cm}
\textbf{[IMAGE : Comparaison ex\'{e}cution s\'{e}quentielle (CPU) vs.\ parall\`{e}le (FPGA)]}\\[0.3cm]
\small CPU : $\sim$50 instructions s\'{e}quentielles par ronde.\\
FPGA : toutes les op\'{e}rations d'une ronde en 1 seul cycle d'horloge.
\vspace{3cm}}}
\caption{Parall\'{e}lisme : 50 instructions CPU vs.\ 1 cycle FPGA par ronde}
\label{fig:parallel_comparison}
\end{figure}

\subsection{Raison 2 : Op\'{e}rations \`{a} co\^{u}t nul en mat\'{e}riel}

Certaines op\'{e}rations qui co\^{u}tent des cycles CPU sont \textbf{gratuites} en
mat\'{e}riel :
\begin{itemize}
    \item \textbf{Rotations (ROTR)} : simple reconnexion de fils, co\^{u}t z\'{e}ro
          en surface et en temps ;
    \item \textbf{XOR bit-\`{a}-bit} : une seule couche de portes logiques
          ($<$ 1 ns de d\'{e}lai) ;
    \item \textbf{S\'{e}lection de $K_t$} : multiplexeur c\^{a}bl\'{e}, pas de lecture
          m\'{e}moire.
\end{itemize}

En logiciel, chaque rotation n\'{e}cessite 2--3 instructions (shift + OR + mask).

\subsection{Raison 3 : Pas d'overhead d'instruction}

Le processeur g\'{e}n\'{e}raliste passe une grande partie de son temps \`{a} :
\begin{itemize}
    \item \textbf{Fetch} : lire l'instruction depuis la m\'{e}moire ;
    \item \textbf{Decode} : interpr\'{e}ter l'opcode ;
    \item \textbf{Acc\`{e}s registres/cache} : lire les op\'{e}randes.
\end{itemize}

Notre acc\'{e}l\'{e}rateur n'a \textbf{aucun} de ces overheads. Chaque cycle d'horloge
est enti\`{e}rement productif.

\subsection{Comparaison quantitative}

Le tableau~\ref{tab:perf_comparison} compare les performances de notre
acc\'{e}l\'{e}rateur avec diff\'{e}rentes plateformes logicielles.

\begin{table}[H]
\centering
\caption{Comparaison de performance : logiciel vs.\ acc\'{e}l\'{e}rateur mat\'{e}riel}
\label{tab:perf_comparison}
\begin{tabular}{l r r r r}
\toprule
\textbf{Plateforme} & \textbf{Cyc/bloc} & \textbf{Fr\'{e}q.} & \textbf{D\'{e}bit} & \textbf{Gain} \\
\midrule
ARM Cortex-M4 (logiciel)       & 3\,500 & 168 MHz & 25 Mb/s   & 1$\times$ \\
ARM Cortex-A53 (logiciel)      & 1\,800 & 1.2 GHz & 341 Mb/s  & 14$\times$ \\
Intel i7 (logiciel + AVX2)     & 1\,000 & 4 GHz   & 2\,048 Mb/s & 82$\times$ \\
\midrule
\textbf{Notre acc\'{e}l\'{e}rateur} & \textbf{85} & 100 MHz & \textbf{602 Mb/s} & \textbf{24$\times$} \\
Notre acc\'{e}l\'{e}rateur (200 MHz) & 85 & 200 MHz & 1\,204 Mb/s & 48$\times$ \\
\bottomrule
\end{tabular}
\end{table}

\textbf{Interpr\'{e}tation} : \`{a} seulement 100 MHz, notre acc\'{e}l\'{e}rateur FPGA
surpasse un ARM Cortex-M4 cadenc\'{e} \`{a} 168 MHz d'un facteur \textbf{24}. \`{A}
200 MHz (r\'{e}aliste sur Xilinx Artix-7), le facteur atteint \textbf{48}. De
plus, pendant que le co-processeur calcule, le CPU h\^{o}te est libre pour
d'autres t\^{a}ches.

\subsection{D\'{e}bit th\'{e}orique}

Pour un bloc de 512 bits trait\'{e} en 85 cycles :
\[
\text{D\'{e}bit} = \frac{512 \text{ bits}}{85 \times 10 \text{ ns}} = 602{,}4 \text{ Mbit/s}
\]

Pour un message multi-blocs de $N$ blocs (d\'{e}bit asymptotique) :
\[
\text{Cycles} \approx N \times 67 + 18 \quad\Rightarrow\quad
\text{D\'{e}bit}_{N\to\infty} = \frac{512}{67 \times 10\text{ ns}} = 764 \text{ Mbit/s}
\]

\subsection{Sch\'{e}ma r\'{e}capitulatif de l'acc\'{e}l\'{e}ration}

La figure~\ref{fig:accel_summary} r\'{e}sume visuellement pourquoi notre design
constitue un acc\'{e}l\'{e}rateur.

\begin{figure}[H]
\centering
% =====================================================================
% IMAGE PLACEHOLDER: Sch\'{e}ma r\'{e}capitulatif de l'acc\'{e}l\'{e}ration
% Ins\'{e}rer ici un diagramme montrant :
%   - \`{A} gauche : CPU traitant SHA-256 en ~3500 cycles
%   - \`{A} droite : Co-processeur en ~85 cycles
%   - Fl\`{e}che montrant le gain x41
%   - Annotations : "CPU lib\'{e}r\'{e} pour autres t\^{a}ches"
% =====================================================================
\fbox{\parbox{0.85\textwidth}{\centering\vspace{3cm}
\textbf{[IMAGE : R\'{e}capitulatif -- CPU seul vs.\ CPU + Acc\'{e}l\'{e}rateur]}\\[0.3cm]
\small Gauche : CPU seul (3500 cycles, bloqu\'{e}).\\
Droite : CPU + co-processeur (85 cycles, CPU libre en parall\`{e}le).\\
Gain : $\sim$41$\times$ en cycles, CPU disponible pour autres t\^{a}ches.
\vspace{3cm}}}
\caption{Synth\`{e}se : le co-processeur acc\'{e}l\`{e}re le calcul et lib\`{e}re le CPU}
\label{fig:accel_summary}
\end{figure}


% =============================================================================
\section{Probl\`{e}mes rencontr\'{e}s et corrections}
\label{sec:bugs}
% =============================================================================

Lors de la mise au point par simulation, trois anomalies ont \'{e}t\'{e} identifi\'{e}es
et corrig\'{e}es. Chaque bug est document\'{e} avec son sympt\^{o}me, sa cause racine, et
la correction appliqu\'{e}e.

\subsection{Bug \#1 : D\'{e}calage d'un mot dans la lecture m\'{e}moire}

\begin{description}[leftmargin=!,labelwidth=2cm]
\item[Sympt\^{o}me] Pour les messages non vides, tous les mots $W_i$ \'{e}taient
    d\'{e}cal\'{e}s d'une position : le premier mot lu \'{e}tait le mot de contr\^{o}le
    (nombre de blocs) au lieu de $W_0$.
\item[Cause] La latence d'un cycle de la RAM synchrone n'\'{e}tait pas
    correctement compens\'{e}e. L'adresse \texttt{0x0004} (correspondant \`{a} $W_0$)
    \'{e}tait positionn\'{e}e trop tard sur le bus d'adresses.
\item[Correction] Avancer l'\'{e}mission de l'adresse $W_0$ d'un cycle en la
    pla\c{c}ant d\`{e}s la transition IDLE$\rightarrow$READ\_NBLK, et adapter
    l'arithm\'{e}tique d'adresses en LOAD (\texttt{counter+2} au lieu de
    \texttt{counter+1}).
\end{description}

\subsection{Bug \#2 : Mot $H_0$ incorrect dans le condens\'{e}}

\begin{description}[leftmargin=!,labelwidth=2cm]
\item[Sympt\^{o}me] Pour la cha\^{i}ne vide, $H_1 \ldots H_7$ \'{e}taient corrects
    mais $H_0$ restait \`{a} la valeur initiale \texttt{0x6A09E667}.
\item[Cause] L'accumulation finale ($H_i = H_i + a_i$) et l'\'{e}criture en
    m\'{e}moire se faisaient dans le m\^{e}me cycle. Le premier mot \'{e}crit utilisait
    la valeur de \texttt{H\_result} \emph{avant} sa mise \`{a} jour.
\item[Correction] Insertion d'un \'{e}tat \textbf{ACCUM} s\'{e}par\'{e} entre COMPUTE et
    WRITEBACK, laissant un cycle complet pour la stabilisation du registre.
\end{description}

\subsection{Bug \#3 : D\'{e}bordement d'index sur les constantes $K$}

\begin{description}[leftmargin=!,labelwidth=2cm]
\item[Sympt\^{o}me] Erreur d'ex\'{e}cution en simulation : \texttt{index (64) out of
    bounds (0 to 63)}.
\item[Cause] Le signal \texttt{counter\_prev} pouvait atteindre 64 (valeur de
    \texttt{counter} au dernier cycle de COMPUTE), mais le tableau
    \texttt{K\_cons} n'est index\'{e} que de 0 \`{a} 63.
\item[Correction] Ajout d'une condition de garde \texttt{counter\_prev < 64}
    dans l'expression combinatoire de s\'{e}lection de $K_t$.
\end{description}

\subsection{Le\c{c}ons apprises}

Ces bugs illustrent deux difficult\'{e}s classiques de la conception num\'{e}rique :
\begin{enumerate}
    \item La gestion rigoureuse du \textbf{pipeline d'adresses} lors de
          l'interfaçage avec des m\'{e}moires synchrones \`{a} latence ;
    \item L'importance de v\'{e}rifier les \textbf{conditions aux limites} (premier
          et dernier cycle d'un \'{e}tat) par simulation exhaustive.
\end{enumerate}

% =============================================================================
\section{Conclusion et perspectives}
\label{sec:conclu}
% =============================================================================

\subsection{Bilan}

Ce projet a permis de concevoir et valider un \textbf{acc\'{e}l\'{e}rateur mat\'{e}riel
SHA-256} complet et fonctionnel :
\begin{itemize}
    \item Conforme \`{a} la norme FIPS 180-4 ;
    \item Support multi-blocs (messages de longueur arbitraire) ;
    \item Valid\'{e} avec les vecteurs NIST incluant le cas 2 blocs ;
    \item D\'{e}bit th\'{e}orique de \textbf{602 Mbit/s} \`{a} 100 MHz ;
    \item Gain de performance de \textbf{24 \`{a} 48 fois} par rapport \`{a} un
          processeur embarqu\'{e} (ARM Cortex-M4).
\end{itemize}

Le terme \og \textbf{acc\'{e}l\'{e}rateur} \fg{} est pleinement justifi\'{e} : le circuit
exploite le parall\'{e}lisme spatial, les op\'{e}rations \`{a} co\^{u}t nul (rotations),
et l'absence d'overhead d'instruction pour atteindre des performances
inaccessibles en logiciel sur un processeur embarqu\'{e}.

\subsection{Perspectives d'am\'{e}lioration}

\begin{enumerate}
    \item \textbf{Interface AXI-Lite} : encapsulation dans un wrapper bus
          standard pour int\'{e}gration dans un SoC Xilinx Zynq ;
    \item \textbf{Pr\'{e}-calcul $K_t + W_t$} : registrer la somme un cycle \`{a}
          l'avance pour r\'{e}duire le chemin critique de $\sim$25\% et
          augmenter la Fmax ;
    \item \textbf{Architecture multi-c\oe{}urs} : instancier $N$ c\oe{}urs SHA-256
          en parall\`{e}le pour multiplier le d\'{e}bit par $N$ (applications
          blockchain/mining) ;
    \item \textbf{Support HMAC-SHA256} : ajouter la couche d'authentification
          de messages ;
    \item \textbf{Synth\`{e}se FPGA} : caract\'{e}riser les ressources (LUT, FF, BRAM)
          et mesurer la Fmax r\'{e}elle sur Artix-7 / Zynq-7000.
\end{enumerate}

% =============================================================================
% R\'{e}f\'{e}rences bibliographiques
% =============================================================================
\newpage
\begin{thebibliography}{9}

\bibitem{fips180}
National Institute of Standards and Technology (NIST),
\emph{Secure Hash Standard (SHS)},
FIPS PUB 180-4, ao\^{u}t 2015.
Disponible : \url{https://doi.org/10.6028/NIST.FIPS.180-4}

\bibitem{mcloone}
M.~McLoone et J.~McCanny,
``Efficient Single-Chip Implementation of SHA-384 and SHA-512'',
\emph{Proceedings of IEEE FPL}, 2002.

\bibitem{dadda}
L.~Dadda, M.~Macchetti et J.~Owen,
``An ASIC Design for a High Speed Implementation of the Hash Function SHA-256'',
\emph{Proceedings of GLSVLSI}, 2004.

\bibitem{icact2019}
M.~Dowenbergh et al.,
``Acceleration of the Secure Hash Algorithm-256 (SHA-256) on an FPGA-CPU Cluster
Using OpenCL'',
\emph{Proceedings of ICACT}, 2019.

\bibitem{unfolding}
B.~Padma Priya et al.,
``FPGA-based Implementation of SHA-256 with Improvement of Throughput using
Unfolding Transformation'',
\emph{IEEE Access}, 2022.

\end{thebibliography}

\end{document}
